\section{Conclusion}

\label{2305:sec:conclusion}
In this work we have derived a sharp formula for how many samples are required for a two-layer neural network trained with one-pass SGD to learn a quadratic target in the high-dimensional limit. In particular, we have shown that increasing the width can only improve the sample complexity by a pre-factor, with the overall scaling with the dimension remaining of the same order $n=O(d\log{d})$. Therefore, for this target overparameterization does not significantly help optimization, providing a prototypical example of a hard class of functions for learning with SGD. Our results rely on a low-dimensional description of SGD in terms of a stochastic process describing the evolution of the sufficient statistics in the high-dimensional limit. Surprisingly, we have shown that deriving the sample complexities from the deterministic drift of this process with small initial correlation with the target provides a fairly good approximation of the exit time as computed from the full process with zero initial correlation, showing stochasticity does not play a crucial role in escaping mediocrity for this problem.
% %%%%%%%%%%%%%%%%%%%%%%%%%%%%%%%%%%%%%%%%%%%%%%%%%%%%%%%%%%%%%%%%%%%%%%%%%%%%%%%
